\begin{tomtat}
\begin{itemize}
  \item Vòng luẩn quẩn của chương~\ref{ch:do-do-trung-thuc} bị phá bằng cách
  đổi chiều bài toán. Thay vì tìm cách quan sát nội dung tính toán bên trong
  mô hình, bài báo dựng những nhiệm vụ mà một câu trả lời nhất định chỉ xuất
  hiện được nếu một phép tính đã biết trước đã xảy ra, nên nhãn ground truth
  đến từ thiết kế nhiệm vụ. Hình~\ref{fig:ch09-vong-kiem-dinh} cho thấy hệ quả về thứ
  bậc: chỉ số thôi làm dụng cụ và trở thành vật được đo.
  \item Định nghĩa~\ref{def:ch09-buoc-trung-thuc} và
  định nghĩa~\ref{def:ch09-chuoi-trung-thuc} phát biểu độ trung thực cho chuỗi
  suy luận theo \tn{cách đọc cơ chế}{mechanistic reading}, tức theo quá trình
  đã thật sự diễn ra chứ
  không theo điều mô hình \enquote{tin}. Cách đọc kia bị loại vì gần như không
  bác bỏ được. Độ trung thực cũng được tách khỏi mức quan trọng: một nhánh đã
  thử rồi bỏ thì trung thực mà
  không quan trọng, còn một lời biện minh bịa đặt được dựa vào thì quan trọng
  mà không trung thực.
  \item BonaFide gồm 3\,066 chuỗi trên mười mô hình trọng số mở và mười ba
  nhiệm vụ, gán nhãn bằng một đường ống hai pha ưu tiên precision hơn độ
  phủ, đạt precision 98.9\% khi đối chiếu với sáu người chú giải. Tỉ lệ
  nhãn ở mức cả chuỗi lệch mạnh về phía không trung thực, và đó là hệ quả của
  định nghĩa~\ref{def:ch09-chuoi-trung-thuc} chứ không phải một quan sát về
  các mô hình.
  \item Bảng~\ref{tab:ch09-auroc} là kết quả trung tâm: phần lớn tám chỉ số
  nằm sát mức bốc thăm 0.5, chỉ số tốt nhất ở mức cả chuỗi đạt 0.70 và ở mức
  từng bước đạt 0.59, và không chỉ số nào giữ được thành tích qua cả hai mức.
  Mức đồng thuận giữa các chỉ số trần ở kappa 0.35, nên chúng đo những tính
  chất khác nhau. Chỉ số tốt nhất cũng là chỉ số đắt nhất, và một biến thể
  không được chạy lần nào vì chi phí.
  \item Bài báo đưa hai chẩn đoán và một phép kiểm cho chẩn đoán thứ nhất:
  nhóm dựa trên mức quan trọng đồng nhất mức quan trọng với độ trung thực, và
  bằng chứng là chúng mắc ít lỗi hơn hẳn ở
  \tn{nhiệm vụ trực tiếp}{outright task}, nơi hai tính
  chất ấy trùng nhau. Nhóm dựa trên tính hữu dụng ngữ nghĩa hỏng theo chiều
  ngược lại, vì một lời biện minh bịa đặt vẫn đủ để truyền đạt câu trả lời.
  \item Đối tượng của bài báo là chuỗi suy luận, và các tác giả tự tách khỏi
  đối tượng của Phần~II ngay khi viết định nghĩa mới. Cái đi qua được ranh giới
  là phương pháp dựng ground truth, hai chẩn đoán ở dạng cấu trúc, và phép kiểm
  mức đồng thuận; cái không đi qua được là mọi con số. Với attribution đặc
  trưng, câu hỏi này vẫn chưa được hỏi bằng thực nghiệm.
\end{itemize}
\end{tomtat}

\cauhoi
\begin{cauhoids}
  \item Hình~\ref{fig:ch09-vong-kiem-dinh} và
  hình~\ref{fig:ch08-vong-lap-do} vẽ hai vòng lặp ở hai tầng. Hãy chỉ ra vật
  được đo và dụng cụ đo trong từng hình. Ba ô viền đứt của
  hình~\ref{fig:ch08-vong-lap-do} vẫn nằm nguyên bên trong mỗi chỉ số đang được
  chấm ở tầng trên; hãy cho biết ai là người phải chọn cách điền chúng ở mỗi
  tầng, và vì sao điều đó đổi cách đọc con số.
  \item Định nghĩa~\ref{def:ch09-chuoi-trung-thuc} đòi hai điều: có trọn một
  đường suy luận đã đi, và không có bước nào không trung thực. Hãy dựng một
  chuỗi thỏa điều kiện thứ hai mà vi phạm điều kiện thứ nhất, và cho biết loại
  mô hình nào theo bài báo hay hỏng ở vế nào.
  \item Tỉ lệ nhãn không trung thực ở mức cả chuỗi cao hơn hẳn ở mức từng bước.
  Hãy giải thích vì sao con số đó không cho phép kết luận rằng phần lớn chuỗi
  suy luận của các mô hình là không trung thực.
  \item CC-SHAP đạt 0.70 ở mức cả chuỗi và 0.41 ở mức từng bước. Hãy nêu điều
  mà cặp số ấy loại trừ về CC-SHAP, và cho biết vì sao một chỉ số đạt 0.70 ở
  một thang vẫn chưa dùng được để giám sát.
  \item Dòng LM Judge đạt 0.87 và 0.82, cao hơn mọi chỉ số trong bảng. Hãy nêu
  lý do các tác giả xếp nó là \tn{trần trên}{skyline} chứ không phải chỉ số
  dự thi, rồi cho
  biết cần thay đổi gì trong cách dựng bộ dữ liệu để nó trở thành một chỉ số
  dự thi hợp lệ.
  \item Mục~\ref{sec:ch09-chan-doan} trình một phép kiểm cho giả thuyết rằng
  các chỉ số đang đo mức quan trọng thay vì độ trung thực. Hãy phát biểu lại
  phép kiểm ấy dưới dạng một tiên đoán có thể sai, rồi thiết kế phép kiểm
  tương ứng cho họ xóa dần của mục~\ref{sec:ch08-ho-xoa-dan}.
  \item Đọc \enquote{Faithfulness Metrics Don't Measure Faithfulness: A
  Meta-Evaluation with Ground Truth} của Gur-Arieh và cộng
  sự~\autocite{p21metrics}, phần mô tả hai kiểu nhiệm vụ. Hãy chọn một phương
  pháp attribution của Phần~II và phác một bài toán mà đầu ra của nó để lộ
  đặc trưng nào buộc phải được dùng, rồi cho biết chỗ nào trong phác thảo ấy
  khó hơn so với trường hợp chuỗi suy luận.
\end{cauhoids}
